\scanpage{7}

\deleted{En d'autres termes}\inserted{\unsure{i.e.|r.e.|r.c.}}, si \inserted{lorsque $Y$ est \inserted{\unsure{déjà|supposé}} fermé dans $X$} on a $F(T_0)\ne\emptyset$, $F(U)\ne\emptyset$, alors $\deleted{\unsure{x|X}}$ $F(T)\ne\emptyset$. \inserted{En fait, il suffit ici $F(U)\ne\emptyset$ pour $F(T)\ne\unsure{0|\emptyset}$.}
\deleted{\unsure{D} si $X_{T_0}=Y_{T_0}$ et} \inserted{En d'autres termes si} $X_{\deleted{\unsure{T}}U}=Y_U$, alors $X=Y$.
En effet, comme $\deleted{\unsure{S}}T$ est l'adhérence schématique de $U$ et $X_T$ plat sur $T$, $X_{\deleted{\unsure{S}}T}$ est l'adhérence schématique de $X_U$, donc \deleted{on gagne} $Y_T=X_T$ puisque $Y_T$ sous-préschéma fermé de $X_T$ qui \unsure{majore|possède|majorant} $X_U$.

$8^\circ)$ Si $T$ sur $S$ irréductible muni de $J$ de carré nul, $T_0=V(J)$, si $F(T_0)\ne\emptyset$, \deleted{alors $\exists$ ouvert dense \unsure{tel|de} $F(\underline{\mathcal O}_{T,\unsure{y|\eta}})\ne\emptyset$} alors $\exists$ ouvert dense $U$ de $T$ \deleted{sur lequel $F$} tel que $F_U$ repr.

% slanted
\margin{$X$ plat t.f.}
En effet, il suffit que $\deleted{\unsure{F|X}}X_U$ soit ``ess. libre'' sur $U$, et on sait qu'on peut trouver $U$ de cette façon EGA IV 11

\deleted{\unsure{D}} De plus, on a vérifié $7^\circ$ sous la forme plus forte \unsure{signifiant} que le foncteur est \uuline{propre}.

\deleted{\uline{Corollaire} \unsure{So}}

\deleted{\uline{Corollaire} Supposons que $X$ soit un $S$-groupe, $Y$ un $S$ sous-groupe. Supposons que $X$ ait la propriété suivante :

Pour tout \deleted{\unsure{sch}} $T$ \inserted{sur $S$} spectr. d'un anneau de valuation discrète, \deleted{sur $S$,} et tout \deleted{sg \unsure{pr}} sous groupe \deleted{\unsure{fermé}} $\unsure{Z|H}$ de $X_T$ qui \deleted{\unsure{contienne|contient} la fibre fermée \unsure{entière} spéciale} \unsure{majore en|ou pour un}}
% archive sheet number: 6
